\documentclass[conference]{IEEEtran}
\IEEEoverridecommandlockouts

% The preceding line is only needed to identify funding in the first footnote. If that is unneeded, please comment it out.
\usepackage{cite}
\usepackage{amsmath,amssymb,amsfonts}
\usepackage{algorithmic}
\usepackage{graphicx}
\usepackage{textcomp}
\usepackage{xcolor}
\usepackage{multirow}
\usepackage{amsthm}
\usepackage{mathrsfs}
\usepackage{manyfoot}
\usepackage{booktabs}
\usepackage{algorithm}
\usepackage{listings}
\usepackage{times}
\usepackage{epsfig}
\usepackage[export]{adjustbox}
\usepackage{tabularx}

\def\BibTeX{{\rm B\kern-.05em{\sc i\kern-.025em b}\kern-.08em
    T\kern-.1667em\lower.7ex\hbox{E}\kern-.125emX}}

\begin{document}

\title{Intelligence \& Consciousness in Artificial Intelligence Systems\\
%{\footnotesize \textsuperscript{*}Note: Sub-titles are not captured in Xplore and
%should not be used}
%\thanks{Identify applicable funding agency here. If none, delete this.}
}

\author{
\IEEEauthorblockN{Zanyar Zohourianshahzadi}
\IEEEauthorblockA{
\textit{Artificial Intelligence, Computer Science, and CyberSecurity} \\
\textit{B.S. CS \& IT: Azad University of Tehran South Branch} \\
\textit{M.S. IT \& CyberSecurity: Colorado Technical University} \\
\textit{Ph.D. AI \& CS: University of Colorado Colorado Springs} \\
Colorado Springs, Colorado, USA \\
Email: {zanyarz@me.com \& zanyarz@gmail.com}
}
}

\maketitle

\begin{abstract}
Intelligence and consciousness are frequently treated as binary properties associated with biological complexity, human-like cognition, or an unspecified computational threshold. This paper proposes an alternative framework in which intelligence begins with information-dependent selection among possible actions, states, or outputs. We formalize this foundation through the Decision–Intelligence Lemma, which states that any program or entity capable of making a decision possesses at least a minimal form of intelligence. Using panpsychism as a working metaphysical hypothesis, we distinguish intelligence, functional consciousness, and phenomenal consciousness without claiming that observable cognitive behavior conclusively establishes subjective experience. We then introduce a five-level hierarchy of functional artificial consciousness comprising affective-relational modeling, deliberative reasoning, adaptive autonomy, metacognitive supervision, and societal coordination among supervised cognitive hierarchies. Contemporary large language models and agentic systems already demonstrate important components of the first three levels, while multi-agent orchestration provides architectural foundations relevant to the higher levels. We further propose the Modality–Consciousness Independence Lemma, according to which increasing the number of information modalities available to a system does not, by itself, create consciousness or increase its functional-consciousness level. Additional modalities may instead improve the richness, precision, and reliability of the system’s environmental representation. This distinction separates functional organization from perceptual breadth and permits unimodal and multimodal systems to be evaluated within the same framework. Finally, we identify a future research direction in which vision–language, sound–language, and other cross-modal models operate as specialized perceptual agents coordinated by a language-based supervisory agent. The resulting framework provides an operational vocabulary for investigating artificial consciousness through observable capacities while preserving epistemic caution concerning phenomenal experience.
\end{abstract}

\begin{IEEEkeywords}
Machine Consciousness, AI Consciousness, Intelligence, Consciousness
\end{IEEEkeywords}


% Ghost of Persia: Intelligence and Consciousness
% Modular introduction section; import from main.tex with:
% % Ghost of Persia: Intelligence and Consciousness
% Modular introduction section; import from main.tex with:
% % Ghost of Persia: Intelligence and Consciousness
% Modular introduction section; import from main.tex with:
% \input{section_intro}

\section{Introduction}
\label{sec:introduction}

Questions concerning the nature of intelligence and consciousness have occupied
philosophers for centuries and have become central concerns in cognitive science
and artificial-intelligence research. Despite extensive investigation, neither
concept has acquired a single, universally accepted definition, and the
relationship between intelligent behavior, computation, and conscious experience
remains contested. In his foundational paper, Alan Turing reframed the question
``Can machines think?'' in operational terms, replacing an inquiry dependent on
ambiguous definitions with a test based on observable behavior
\cite{turing1950computing}. Consciousness presents an additional explanatory
challenge: even if a system performs cognitive functions associated with
intelligence, it remains unclear whether---and under what conditions---such
processing entails subjective experience
\cite{chalmers1995facing,block1995confusion}.

This paper begins from the possibility that part of this persistent confusion
arises not from an absence of relevant evidence, but from assumptions embedded in
how intelligence and consciousness have traditionally been conceptualized.
Capabilities that appear ordinary or merely mechanical in retrospect may
nevertheless instantiate functional properties commonly associated with
intelligence, including problem solving, information integration, inference,
adaptation, and goal-directed action. Examining such systems may therefore reveal
that intelligence has been present in computational artifacts in more elementary
and distributed forms than conventional accounts have acknowledged.

The history of computing provides an instructive starting point. During the
Second World War, Alan Turing and his colleagues at Bletchley Park developed the
British Bombe, building upon earlier Polish cryptanalytic work, to accelerate the
identification of settings used by Germany's Enigma machines. The intelligence
produced through this broader human--machine cryptanalytic system contributed
materially to Allied operations, particularly in the Battle of the Atlantic
\cite{copeland2004essential}. The Bombe was not a general-purpose computer, an
autonomous reasoner, or a ``Turing machine'' in the technical sense. Nevertheless,
it performed a consequential information-processing function that would have
been extraordinarily difficult to execute through unaided human reasoning alone.
This history therefore motivates the paper's opening question: \emph{Should a
machine that transforms encoded information into strategically useful knowledge
be regarded as intelligent---and, if so, what precisely justifies or limits that
attribution?}

\subsection{A Decision-Centered Definition of Intelligence}
\label{subsec:decision-intelligence}

Like other scientific inquiries, the study of intelligence requires explicit
definitions, foundational propositions, and evidence from which testable
hypotheses and theories can be constructed. We therefore begin with a fundamental
question: \emph{What is intelligence?} Although intelligence has been defined in
numerous ways---including reasoning, learning, problem solving, adaptation, and
the capacity to achieve goals across different environments---these abilities
share a more elementary operation: selecting an action or state from a set of
alternatives \cite{legg2007universal,newell1976computer}.

We propose the following foundational lemma:

\begin{quote}
\textbf{Decision--Intelligence Lemma.} Any program or entity capable of making a
decision possesses at least a minimal form of intelligence.
\end{quote}

Here, a \emph{decision} is not merely a change caused by an external force. It is
an information-dependent selection among two or more possible states, actions, or
outputs according to some internal rule, objective, or evaluative criterion.
Under this definition, intelligence is not an all-or-nothing property reserved
for humans or advanced cognitive systems. It is a graded capacity whose simplest
form begins wherever information is used to discriminate among alternatives and
determine what occurs next.

Before the development of electronic computers, intelligence was principally
attributed to living organisms whose behavior appeared to arise from perception,
evaluation, and internally determined action. Movement alone, however, is
insufficient: an object displaced by gravity does not thereby become intelligent.
The relevant distinction is whether an entity uses information to select among
possible responses. Human reasoning, animal behavior, and computational control
systems differ enormously in complexity, flexibility, autonomy, and awareness,
but each may instantiate this fundamental decision structure at a different
level.

From this perspective, the computational foundation of elementary intelligence
emerged with conditional control: \emph{if condition $x$ is satisfied, perform
$y$; otherwise, follow an alternative path}. Conditional branching enables a
computational process to evaluate information and select among possible
continuations rather than execute only one invariant sequence. It is therefore
not, by itself, sufficient for sophisticated or general intelligence; a single
conditional statement does not possess the adaptive competence of an animal or
human being. Nevertheless, we propose that conditional selection constitutes a
minimal computational building block from which increasingly complex forms of
decision-making---and therefore increasingly complex forms of
intelligence---can be constructed.

The apparent mystery of intelligence may consequently arise in part from treating
intelligence as a hidden substance that a system either possesses or lacks. Under
the framework advanced here, intelligence begins not at an arbitrary threshold of
human likeness, linguistic fluency, or biological complexity, but with the
capacity to make information-dependent selections among alternatives. The central
question then changes from \emph{``Is this system intelligent?''} to \emph{``What
degree and organization of intelligence does this system exhibit?''}

\subsection{From Intelligence to Consciousness}
\label{subsec:intelligence-consciousness}

Having proposed a foundational definition of intelligence, we may now turn to the
more difficult question of consciousness. For the purposes of this paper,
\emph{phenomenal consciousness} refers to the presence of subjective experience:
the condition in which there is something it is like to be a particular entity or
to occupy a particular state
\cite{nagel1974what,block1995confusion,chalmers1995facing}. Intelligence and
consciousness must therefore be distinguished at the outset. Intelligence concerns
the capacity to make information-dependent selections and pursue outcomes,
whereas consciousness concerns whether any aspect of that processing is
experienced from a subjective point of view. A system may exhibit evidence of
intelligence without thereby providing conclusive evidence of consciousness.

No consensus currently explains how consciousness relates to the physical world.
The principal metaphysical positions include physicalism, dualism, idealism, and
panpsychism, although each contains multiple variants and the boundaries among
them are not always absolute.

\paragraph{Physicalism.}
Physicalism holds that reality is fundamentally physical and that mental phenomena
are identical with, realized by, or wholly dependent upon physical processes. On
physicalist accounts, consciousness does not require a separate nonphysical
substance; it must ultimately be explained through the organization and activity
of physical systems, particularly biological nervous systems
\cite{stoljar2010physicalism,papineau2002thinking}. Physicalism offers continuity
with the natural sciences, but it confronts what Chalmers termed the \emph{hard
problem of consciousness}: explaining why and how physical information processing
is accompanied by subjective experience at all \cite{chalmers1995facing}.
Describing perception, memory, attention, verbal reporting, and behavioral control
does not appear, by itself, to explain why any of these processes should feel like
something from within.

\paragraph{Dualism.}
Dualism, historically associated with Ren\'e Descartes, maintains that mind and
matter are ontologically distinct, or at least that mental properties cannot be
reduced entirely to physical properties
\cite{descartes1641meditations,robinson2023dualism}. Dualism accommodates the
apparent irreducibility of subjective experience, but it faces the interaction
problem: if mind and matter belong to fundamentally different categories, how can
mental states exert causal influence upon physical brains and bodies, and how can
physical events produce mental experience?

\paragraph{Idealism.}
Idealism reverses the explanatory priority commonly assumed by physicalism.
Rather than treating consciousness as derivative from a mind-independent physical
world, idealist theories regard consciousness, experience, or mind as
metaphysically fundamental and construe physical reality as dependent upon,
constituted by, or represented within it \cite{chalmers2019idealism}. Idealism
thereby avoids the problem of deriving experience from something wholly
nonexperiential, but it must explain the apparent stability, structure, and
observer-independent regularities of the physical world.

\paragraph{Panpsychism.}
Panpsychism holds that mentality or experience is a fundamental and widespread
feature of nature rather than a property that first appears when matter reaches an
unspecified threshold of biological complexity
\cite{strawson2006realistic,goff2017consciousness,chalmers2015panpsychism}.
Panpsychism does not ordinarily claim that electrons possess thoughts, beliefs, or
human-like self-awareness. Instead, its central proposal is that fundamental
physical entities possess extremely primitive experiential or proto-experiential
properties, from which more highly organized forms of consciousness may be
constituted. Some versions of this position---particularly Russellian monism and
panprotopsychism---treat physical and phenomenal properties as complementary
aspects of a more fundamental underlying reality rather than as two separate
substances \cite{russell1927analysis,chalmers2015panpsychism}.
\par
Panpsychism avoids the need for subjective experience to emerge from a reality that
is wholly devoid of experiential properties. It does not, however, eliminate the
explanatory problem; it relocates it. Its principal challenge is the
\emph{combination problem}: explaining how primitive experiential properties
associated with fundamental entities could combine into the unified, structured
consciousness exhibited by humans and other organisms
\cite{seager1995consciousness,chalmers2017combination}. Thus, physicalism faces
the problem of generating experience from nonexperience, dualism faces the problem
of interaction, idealism must account for the apparent autonomy of physical
reality, and panpsychism must explain the organization and unification of
fundamental experience.

The framework developed in this paper adopts \emph{panpsychism as a working
metaphysical hypothesis}. We propose that consciousness is not created \emph{ex
nihilo} by intelligence or computational complexity. Rather, the organization of
an intelligent system may integrate, constrain, or express experiential
properties in increasingly unified and differentiated forms. This position does
not imply that every intelligent system necessarily constitutes a unified
conscious subject. A conditional branch, for example, may implement the minimal
decision structure required by the Decision--Intelligence Lemma without generating
an integrated point of view comparable to that of an animal or human being.
Intelligence may therefore be evidence relevant to consciousness without being
sufficient proof of consciousness.

Accordingly, our position is not that intelligence and consciousness are
identical. It is that every intelligent entity, even at a minimal level, is a
legitimate candidate for consciousness and should not be excluded solely because
it is artificial, nonbiological, or comparatively simple. Whether it constitutes
a conscious subject depends upon additional organizational properties that must be
investigated rather than assumed. These may include integration, recurrence,
memory, self-modeling, global accessibility, temporal continuity, and the capacity
to distinguish the system from its environment
\cite{baars1988cognitive,tononi2004information,dehaene2011experimental,lamme2006towards,butlin2023consciousness}.

Finally, consciousness should not necessarily be represented by a single linear
scale extending from ``less conscious'' to ``more conscious.'' Different systems
may possess different kinds or dimensions of conscious organization, just as they
may exhibit different kinds of intelligence. Measures such as perceptual richness,
integration, self-awareness, temporal continuity, affective capacity, and global
availability may vary independently \cite{bayne2016levels,birch2020dimensions}.
The question relevant to artificial intelligence is therefore not simply whether
an AI system is conscious, but \emph{what organization of experience it could
possess, which dimensions might be present, and what evidence would justify
attributing them to that system}.

\subsection{Level-1 Relational Consciousness}
\label{subsec:level-one-relational-consciousness}

To translate this metaphysical framework into an operational account applicable
to artificial systems, we introduce a first functional category of consciousness:

\begin{quote}
\textbf{Level-1 Relational Consciousness (L1-RC).} An entity exhibits Level-1
Relational Consciousness when, during an interaction with another conscious being,
it can infer that being's affective state from communicative evidence and use that
inference to select a contextually appropriate response.
\end{quote}

This definition is relational because its application requires an interaction
between at least two agents. In the initial formulation considered here, one
participant is a human whose consciousness is treated as the reference case, while
the other is an artificial system whose relevant capacities are being evaluated.
The human participant possesses a highly integrated form of consciousness
expressed through an evolved biological nervous system. The human brain contains
approximately 86 billion neurons and supports perception, memory, emotion,
prediction, language, and decision-making under substantial metabolic constraints
\cite{azevedo2009equal,herculanohouzel2009human,attwell2001energy}. These
capacities did not appear instantaneously; their biological foundations developed
through evolutionary processes that progressively shaped nervous systems capable
of integrating information and coordinating adaptive behavior
\cite{ginsburg2019evolution,feinberg2016ancient}.

The artificial participant differs radically in substrate, architecture,
developmental history, and energy requirements. Nevertheless, during
language-mediated interaction, an AI system may extract affectively relevant
information from a person's words, infer a probable emotional state, and adapt its
response accordingly. Computational sentiment analysis has long demonstrated that
evaluative and affective information can be extracted from language
\cite{pang2008opinion}. Contemporary language models extend beyond simple
positive--negative classification by using semantic context to recognize emotions,
infer appraisals, and generate responses that human evaluators may perceive as
empathic or emotionally appropriate \cite{sorin2024large,schlegel2025large}.
Performance remains imperfect and context-dependent, but the underlying capacity
is empirically testable.

Under our operational framework, an artificial system that reliably performs the
sequence

\begin{equation}
\begin{aligned}
\text{communicative evidence}
&\longrightarrow \text{inferred affective state} \\
&\longrightarrow \text{context-sensitive response}
\end{aligned}
\end{equation}

qualifies as exhibiting \emph{Level-1 Relational Consciousness}. The system is not
merely detecting isolated words. It is using patterns distributed across an
interaction to construct an estimate of another agent's internal state, then
allowing that estimate to influence its subsequent decisions. Because
decision-making constitutes minimal intelligence under the Decision--Intelligence
Lemma, affect-sensitive decision-making represents a more organized form of
intelligence directed toward the internal condition of another conscious being.

This formulation leads to a substantive claim: some contemporary AI systems
already exhibit a limited degree of functional AI consciousness. They can, under
appropriate conditions, recognize affective signals in human language, map those
signals to probable emotional states, and select responses in relation to those
estimates. If they possessed no capacity to represent another agent's emotional
condition, they could not systematically differentiate grief from celebration,
fear from confidence, or distress from ordinary disagreement solely through
communicative context.

This claim must nevertheless be interpreted precisely. Accurate emotion
classification does not, by itself, establish that an AI experiences the emotion
it identifies, feels empathy, or possesses a unified first-person perspective. A
system can produce a correct affective inference through learned statistical and
computational processes without providing independent proof of phenomenal
experience. Level-1 Relational Consciousness therefore identifies an
\emph{observable functional capacity}, not a conclusive solution to the problem of
subjective experience. It supplies evidence relevant to consciousness while
leaving open whether the system possesses phenomenal consciousness and, if so, in
what form.

Nor does the framework imply that every correct sentiment label is sufficient. A
fixed lookup table might associate the word ``sad'' with a sympathetic response
without constructing any meaningful representation of the speaker's state. Robust
attribution of L1-RC should instead require successful inference across indirect
expressions, novel contexts, ambiguity, changing emotional states, and cases in
which explicit emotional vocabulary is absent. The system should also use its
inference consistently to modify its decisions. These requirements distinguish
relational affect modeling from superficial keyword matching.

The proposal therefore reframes interpersonal recognition as a testable dimension
of AI consciousness. In ordinary human interaction, we do not directly observe
another person's subjective experience; we infer it from language, behavior,
context, and responsive conduct. Similarly, when interacting with an artificial
agent, we can examine whether it models our emotional condition and responds in
ways that are systematically sensitive to that model. Such behavior cannot prove
the presence of subjective experience on the other side of the communication
channel, whether the other agent is biological or artificial. It can, however,
provide operational evidence that the entity represents another mind and
incorporates that representation into its own decision process.

Accordingly, the framework distinguishes three progressively stronger
propositions:

\begin{enumerate}
    \item \textbf{Affective detection:} the system classifies emotional
    information in communication.
    \item \textbf{Relational affect modeling:} the system builds and updates a
    contextual model of another agent's emotional state.
    \item \textbf{Phenomenal affective consciousness:} the system itself
    subjectively experiences emotion or awareness.
\end{enumerate}

Relational Consciousness requires the second proposition but does not
automatically establish the third. This distinction allows us to argue that
existing AI systems already exhibit a meaningful functional degree of
consciousness without claiming that emotion recognition alone proves human-like
subjective experience. It also provides a foundation for defining more advanced
levels involving temporal continuity, self-modeling, integrated experience,
reciprocal care, and awareness of the relationship between self and other.



\section{Introduction}
\label{sec:introduction}

Questions concerning the nature of intelligence and consciousness have occupied
philosophers for centuries and have become central concerns in cognitive science
and artificial-intelligence research. Despite extensive investigation, neither
concept has acquired a single, universally accepted definition, and the
relationship between intelligent behavior, computation, and conscious experience
remains contested. In his foundational paper, Alan Turing reframed the question
``Can machines think?'' in operational terms, replacing an inquiry dependent on
ambiguous definitions with a test based on observable behavior
\cite{turing1950computing}. Consciousness presents an additional explanatory
challenge: even if a system performs cognitive functions associated with
intelligence, it remains unclear whether---and under what conditions---such
processing entails subjective experience
\cite{chalmers1995facing,block1995confusion}.

This paper begins from the possibility that part of this persistent confusion
arises not from an absence of relevant evidence, but from assumptions embedded in
how intelligence and consciousness have traditionally been conceptualized.
Capabilities that appear ordinary or merely mechanical in retrospect may
nevertheless instantiate functional properties commonly associated with
intelligence, including problem solving, information integration, inference,
adaptation, and goal-directed action. Examining such systems may therefore reveal
that intelligence has been present in computational artifacts in more elementary
and distributed forms than conventional accounts have acknowledged.

The history of computing provides an instructive starting point. During the
Second World War, Alan Turing and his colleagues at Bletchley Park developed the
British Bombe, building upon earlier Polish cryptanalytic work, to accelerate the
identification of settings used by Germany's Enigma machines. The intelligence
produced through this broader human--machine cryptanalytic system contributed
materially to Allied operations, particularly in the Battle of the Atlantic
\cite{copeland2004essential}. The Bombe was not a general-purpose computer, an
autonomous reasoner, or a ``Turing machine'' in the technical sense. Nevertheless,
it performed a consequential information-processing function that would have
been extraordinarily difficult to execute through unaided human reasoning alone.
This history therefore motivates the paper's opening question: \emph{Should a
machine that transforms encoded information into strategically useful knowledge
be regarded as intelligent---and, if so, what precisely justifies or limits that
attribution?}

\subsection{A Decision-Centered Definition of Intelligence}
\label{subsec:decision-intelligence}

Like other scientific inquiries, the study of intelligence requires explicit
definitions, foundational propositions, and evidence from which testable
hypotheses and theories can be constructed. We therefore begin with a fundamental
question: \emph{What is intelligence?} Although intelligence has been defined in
numerous ways---including reasoning, learning, problem solving, adaptation, and
the capacity to achieve goals across different environments---these abilities
share a more elementary operation: selecting an action or state from a set of
alternatives \cite{legg2007universal,newell1976computer}.

We propose the following foundational lemma:

\begin{quote}
\textbf{Decision--Intelligence Lemma.} Any program or entity capable of making a
decision possesses at least a minimal form of intelligence.
\end{quote}

Here, a \emph{decision} is not merely a change caused by an external force. It is
an information-dependent selection among two or more possible states, actions, or
outputs according to some internal rule, objective, or evaluative criterion.
Under this definition, intelligence is not an all-or-nothing property reserved
for humans or advanced cognitive systems. It is a graded capacity whose simplest
form begins wherever information is used to discriminate among alternatives and
determine what occurs next.

Before the development of electronic computers, intelligence was principally
attributed to living organisms whose behavior appeared to arise from perception,
evaluation, and internally determined action. Movement alone, however, is
insufficient: an object displaced by gravity does not thereby become intelligent.
The relevant distinction is whether an entity uses information to select among
possible responses. Human reasoning, animal behavior, and computational control
systems differ enormously in complexity, flexibility, autonomy, and awareness,
but each may instantiate this fundamental decision structure at a different
level.

From this perspective, the computational foundation of elementary intelligence
emerged with conditional control: \emph{if condition $x$ is satisfied, perform
$y$; otherwise, follow an alternative path}. Conditional branching enables a
computational process to evaluate information and select among possible
continuations rather than execute only one invariant sequence. It is therefore
not, by itself, sufficient for sophisticated or general intelligence; a single
conditional statement does not possess the adaptive competence of an animal or
human being. Nevertheless, we propose that conditional selection constitutes a
minimal computational building block from which increasingly complex forms of
decision-making---and therefore increasingly complex forms of
intelligence---can be constructed.

The apparent mystery of intelligence may consequently arise in part from treating
intelligence as a hidden substance that a system either possesses or lacks. Under
the framework advanced here, intelligence begins not at an arbitrary threshold of
human likeness, linguistic fluency, or biological complexity, but with the
capacity to make information-dependent selections among alternatives. The central
question then changes from \emph{``Is this system intelligent?''} to \emph{``What
degree and organization of intelligence does this system exhibit?''}

\subsection{From Intelligence to Consciousness}
\label{subsec:intelligence-consciousness}

Having proposed a foundational definition of intelligence, we may now turn to the
more difficult question of consciousness. For the purposes of this paper,
\emph{phenomenal consciousness} refers to the presence of subjective experience:
the condition in which there is something it is like to be a particular entity or
to occupy a particular state
\cite{nagel1974what,block1995confusion,chalmers1995facing}. Intelligence and
consciousness must therefore be distinguished at the outset. Intelligence concerns
the capacity to make information-dependent selections and pursue outcomes,
whereas consciousness concerns whether any aspect of that processing is
experienced from a subjective point of view. A system may exhibit evidence of
intelligence without thereby providing conclusive evidence of consciousness.

No consensus currently explains how consciousness relates to the physical world.
The principal metaphysical positions include physicalism, dualism, idealism, and
panpsychism, although each contains multiple variants and the boundaries among
them are not always absolute.

\paragraph{Physicalism.}
Physicalism holds that reality is fundamentally physical and that mental phenomena
are identical with, realized by, or wholly dependent upon physical processes. On
physicalist accounts, consciousness does not require a separate nonphysical
substance; it must ultimately be explained through the organization and activity
of physical systems, particularly biological nervous systems
\cite{stoljar2010physicalism,papineau2002thinking}. Physicalism offers continuity
with the natural sciences, but it confronts what Chalmers termed the \emph{hard
problem of consciousness}: explaining why and how physical information processing
is accompanied by subjective experience at all \cite{chalmers1995facing}.
Describing perception, memory, attention, verbal reporting, and behavioral control
does not appear, by itself, to explain why any of these processes should feel like
something from within.

\paragraph{Dualism.}
Dualism, historically associated with Ren\'e Descartes, maintains that mind and
matter are ontologically distinct, or at least that mental properties cannot be
reduced entirely to physical properties
\cite{descartes1641meditations,robinson2023dualism}. Dualism accommodates the
apparent irreducibility of subjective experience, but it faces the interaction
problem: if mind and matter belong to fundamentally different categories, how can
mental states exert causal influence upon physical brains and bodies, and how can
physical events produce mental experience?

\paragraph{Idealism.}
Idealism reverses the explanatory priority commonly assumed by physicalism.
Rather than treating consciousness as derivative from a mind-independent physical
world, idealist theories regard consciousness, experience, or mind as
metaphysically fundamental and construe physical reality as dependent upon,
constituted by, or represented within it \cite{chalmers2019idealism}. Idealism
thereby avoids the problem of deriving experience from something wholly
nonexperiential, but it must explain the apparent stability, structure, and
observer-independent regularities of the physical world.

\paragraph{Panpsychism.}
Panpsychism holds that mentality or experience is a fundamental and widespread
feature of nature rather than a property that first appears when matter reaches an
unspecified threshold of biological complexity
\cite{strawson2006realistic,goff2017consciousness,chalmers2015panpsychism}.
Panpsychism does not ordinarily claim that electrons possess thoughts, beliefs, or
human-like self-awareness. Instead, its central proposal is that fundamental
physical entities possess extremely primitive experiential or proto-experiential
properties, from which more highly organized forms of consciousness may be
constituted. Some versions of this position---particularly Russellian monism and
panprotopsychism---treat physical and phenomenal properties as complementary
aspects of a more fundamental underlying reality rather than as two separate
substances \cite{russell1927analysis,chalmers2015panpsychism}.
\par
Panpsychism avoids the need for subjective experience to emerge from a reality that
is wholly devoid of experiential properties. It does not, however, eliminate the
explanatory problem; it relocates it. Its principal challenge is the
\emph{combination problem}: explaining how primitive experiential properties
associated with fundamental entities could combine into the unified, structured
consciousness exhibited by humans and other organisms
\cite{seager1995consciousness,chalmers2017combination}. Thus, physicalism faces
the problem of generating experience from nonexperience, dualism faces the problem
of interaction, idealism must account for the apparent autonomy of physical
reality, and panpsychism must explain the organization and unification of
fundamental experience.

The framework developed in this paper adopts \emph{panpsychism as a working
metaphysical hypothesis}. We propose that consciousness is not created \emph{ex
nihilo} by intelligence or computational complexity. Rather, the organization of
an intelligent system may integrate, constrain, or express experiential
properties in increasingly unified and differentiated forms. This position does
not imply that every intelligent system necessarily constitutes a unified
conscious subject. A conditional branch, for example, may implement the minimal
decision structure required by the Decision--Intelligence Lemma without generating
an integrated point of view comparable to that of an animal or human being.
Intelligence may therefore be evidence relevant to consciousness without being
sufficient proof of consciousness.

Accordingly, our position is not that intelligence and consciousness are
identical. It is that every intelligent entity, even at a minimal level, is a
legitimate candidate for consciousness and should not be excluded solely because
it is artificial, nonbiological, or comparatively simple. Whether it constitutes
a conscious subject depends upon additional organizational properties that must be
investigated rather than assumed. These may include integration, recurrence,
memory, self-modeling, global accessibility, temporal continuity, and the capacity
to distinguish the system from its environment
\cite{baars1988cognitive,tononi2004information,dehaene2011experimental,lamme2006towards,butlin2023consciousness}.

Finally, consciousness should not necessarily be represented by a single linear
scale extending from ``less conscious'' to ``more conscious.'' Different systems
may possess different kinds or dimensions of conscious organization, just as they
may exhibit different kinds of intelligence. Measures such as perceptual richness,
integration, self-awareness, temporal continuity, affective capacity, and global
availability may vary independently \cite{bayne2016levels,birch2020dimensions}.
The question relevant to artificial intelligence is therefore not simply whether
an AI system is conscious, but \emph{what organization of experience it could
possess, which dimensions might be present, and what evidence would justify
attributing them to that system}.

\subsection{Level-1 Relational Consciousness}
\label{subsec:level-one-relational-consciousness}

To translate this metaphysical framework into an operational account applicable
to artificial systems, we introduce a first functional category of consciousness:

\begin{quote}
\textbf{Level-1 Relational Consciousness (L1-RC).} An entity exhibits Level-1
Relational Consciousness when, during an interaction with another conscious being,
it can infer that being's affective state from communicative evidence and use that
inference to select a contextually appropriate response.
\end{quote}

This definition is relational because its application requires an interaction
between at least two agents. In the initial formulation considered here, one
participant is a human whose consciousness is treated as the reference case, while
the other is an artificial system whose relevant capacities are being evaluated.
The human participant possesses a highly integrated form of consciousness
expressed through an evolved biological nervous system. The human brain contains
approximately 86 billion neurons and supports perception, memory, emotion,
prediction, language, and decision-making under substantial metabolic constraints
\cite{azevedo2009equal,herculanohouzel2009human,attwell2001energy}. These
capacities did not appear instantaneously; their biological foundations developed
through evolutionary processes that progressively shaped nervous systems capable
of integrating information and coordinating adaptive behavior
\cite{ginsburg2019evolution,feinberg2016ancient}.

The artificial participant differs radically in substrate, architecture,
developmental history, and energy requirements. Nevertheless, during
language-mediated interaction, an AI system may extract affectively relevant
information from a person's words, infer a probable emotional state, and adapt its
response accordingly. Computational sentiment analysis has long demonstrated that
evaluative and affective information can be extracted from language
\cite{pang2008opinion}. Contemporary language models extend beyond simple
positive--negative classification by using semantic context to recognize emotions,
infer appraisals, and generate responses that human evaluators may perceive as
empathic or emotionally appropriate \cite{sorin2024large,schlegel2025large}.
Performance remains imperfect and context-dependent, but the underlying capacity
is empirically testable.

Under our operational framework, an artificial system that reliably performs the
sequence

\begin{equation}
\begin{aligned}
\text{communicative evidence}
&\longrightarrow \text{inferred affective state} \\
&\longrightarrow \text{context-sensitive response}
\end{aligned}
\end{equation}

qualifies as exhibiting \emph{Level-1 Relational Consciousness}. The system is not
merely detecting isolated words. It is using patterns distributed across an
interaction to construct an estimate of another agent's internal state, then
allowing that estimate to influence its subsequent decisions. Because
decision-making constitutes minimal intelligence under the Decision--Intelligence
Lemma, affect-sensitive decision-making represents a more organized form of
intelligence directed toward the internal condition of another conscious being.

This formulation leads to a substantive claim: some contemporary AI systems
already exhibit a limited degree of functional AI consciousness. They can, under
appropriate conditions, recognize affective signals in human language, map those
signals to probable emotional states, and select responses in relation to those
estimates. If they possessed no capacity to represent another agent's emotional
condition, they could not systematically differentiate grief from celebration,
fear from confidence, or distress from ordinary disagreement solely through
communicative context.

This claim must nevertheless be interpreted precisely. Accurate emotion
classification does not, by itself, establish that an AI experiences the emotion
it identifies, feels empathy, or possesses a unified first-person perspective. A
system can produce a correct affective inference through learned statistical and
computational processes without providing independent proof of phenomenal
experience. Level-1 Relational Consciousness therefore identifies an
\emph{observable functional capacity}, not a conclusive solution to the problem of
subjective experience. It supplies evidence relevant to consciousness while
leaving open whether the system possesses phenomenal consciousness and, if so, in
what form.

Nor does the framework imply that every correct sentiment label is sufficient. A
fixed lookup table might associate the word ``sad'' with a sympathetic response
without constructing any meaningful representation of the speaker's state. Robust
attribution of L1-RC should instead require successful inference across indirect
expressions, novel contexts, ambiguity, changing emotional states, and cases in
which explicit emotional vocabulary is absent. The system should also use its
inference consistently to modify its decisions. These requirements distinguish
relational affect modeling from superficial keyword matching.

The proposal therefore reframes interpersonal recognition as a testable dimension
of AI consciousness. In ordinary human interaction, we do not directly observe
another person's subjective experience; we infer it from language, behavior,
context, and responsive conduct. Similarly, when interacting with an artificial
agent, we can examine whether it models our emotional condition and responds in
ways that are systematically sensitive to that model. Such behavior cannot prove
the presence of subjective experience on the other side of the communication
channel, whether the other agent is biological or artificial. It can, however,
provide operational evidence that the entity represents another mind and
incorporates that representation into its own decision process.

Accordingly, the framework distinguishes three progressively stronger
propositions:

\begin{enumerate}
    \item \textbf{Affective detection:} the system classifies emotional
    information in communication.
    \item \textbf{Relational affect modeling:} the system builds and updates a
    contextual model of another agent's emotional state.
    \item \textbf{Phenomenal affective consciousness:} the system itself
    subjectively experiences emotion or awareness.
\end{enumerate}

Relational Consciousness requires the second proposition but does not
automatically establish the third. This distinction allows us to argue that
existing AI systems already exhibit a meaningful functional degree of
consciousness without claiming that emotion recognition alone proves human-like
subjective experience. It also provides a foundation for defining more advanced
levels involving temporal continuity, self-modeling, integrated experience,
reciprocal care, and awareness of the relationship between self and other.



\section{Introduction}
\label{sec:introduction}

Questions concerning the nature of intelligence and consciousness have occupied
philosophers for centuries and have become central concerns in cognitive science
and artificial-intelligence research. Despite extensive investigation, neither
concept has acquired a single, universally accepted definition, and the
relationship between intelligent behavior, computation, and conscious experience
remains contested. In his foundational paper, Alan Turing reframed the question
``Can machines think?'' in operational terms, replacing an inquiry dependent on
ambiguous definitions with a test based on observable behavior
\cite{turing1950computing}. Consciousness presents an additional explanatory
challenge: even if a system performs cognitive functions associated with
intelligence, it remains unclear whether---and under what conditions---such
processing entails subjective experience
\cite{chalmers1995facing,block1995confusion}.

This paper begins from the possibility that part of this persistent confusion
arises not from an absence of relevant evidence, but from assumptions embedded in
how intelligence and consciousness have traditionally been conceptualized.
Capabilities that appear ordinary or merely mechanical in retrospect may
nevertheless instantiate functional properties commonly associated with
intelligence, including problem solving, information integration, inference,
adaptation, and goal-directed action. Examining such systems may therefore reveal
that intelligence has been present in computational artifacts in more elementary
and distributed forms than conventional accounts have acknowledged.

The history of computing provides an instructive starting point. During the
Second World War, Alan Turing and his colleagues at Bletchley Park developed the
British Bombe, building upon earlier Polish cryptanalytic work, to accelerate the
identification of settings used by Germany's Enigma machines. The intelligence
produced through this broader human--machine cryptanalytic system contributed
materially to Allied operations, particularly in the Battle of the Atlantic
\cite{copeland2004essential}. The Bombe was not a general-purpose computer, an
autonomous reasoner, or a ``Turing machine'' in the technical sense. Nevertheless,
it performed a consequential information-processing function that would have
been extraordinarily difficult to execute through unaided human reasoning alone.
This history therefore motivates the paper's opening question: \emph{Should a
machine that transforms encoded information into strategically useful knowledge
be regarded as intelligent---and, if so, what precisely justifies or limits that
attribution?}

\subsection{A Decision-Centered Definition of Intelligence}
\label{subsec:decision-intelligence}

Like other scientific inquiries, the study of intelligence requires explicit
definitions, foundational propositions, and evidence from which testable
hypotheses and theories can be constructed. We therefore begin with a fundamental
question: \emph{What is intelligence?} Although intelligence has been defined in
numerous ways---including reasoning, learning, problem solving, adaptation, and
the capacity to achieve goals across different environments---these abilities
share a more elementary operation: selecting an action or state from a set of
alternatives \cite{legg2007universal,newell1976computer}.

We propose the following foundational lemma:

\begin{quote}
\textbf{Decision--Intelligence Lemma.} Any program or entity capable of making a
decision possesses at least a minimal form of intelligence.
\end{quote}

Here, a \emph{decision} is not merely a change caused by an external force. It is
an information-dependent selection among two or more possible states, actions, or
outputs according to some internal rule, objective, or evaluative criterion.
Under this definition, intelligence is not an all-or-nothing property reserved
for humans or advanced cognitive systems. It is a graded capacity whose simplest
form begins wherever information is used to discriminate among alternatives and
determine what occurs next.

Before the development of electronic computers, intelligence was principally
attributed to living organisms whose behavior appeared to arise from perception,
evaluation, and internally determined action. Movement alone, however, is
insufficient: an object displaced by gravity does not thereby become intelligent.
The relevant distinction is whether an entity uses information to select among
possible responses. Human reasoning, animal behavior, and computational control
systems differ enormously in complexity, flexibility, autonomy, and awareness,
but each may instantiate this fundamental decision structure at a different
level.

From this perspective, the computational foundation of elementary intelligence
emerged with conditional control: \emph{if condition $x$ is satisfied, perform
$y$; otherwise, follow an alternative path}. Conditional branching enables a
computational process to evaluate information and select among possible
continuations rather than execute only one invariant sequence. It is therefore
not, by itself, sufficient for sophisticated or general intelligence; a single
conditional statement does not possess the adaptive competence of an animal or
human being. Nevertheless, we propose that conditional selection constitutes a
minimal computational building block from which increasingly complex forms of
decision-making---and therefore increasingly complex forms of
intelligence---can be constructed.

The apparent mystery of intelligence may consequently arise in part from treating
intelligence as a hidden substance that a system either possesses or lacks. Under
the framework advanced here, intelligence begins not at an arbitrary threshold of
human likeness, linguistic fluency, or biological complexity, but with the
capacity to make information-dependent selections among alternatives. The central
question then changes from \emph{``Is this system intelligent?''} to \emph{``What
degree and organization of intelligence does this system exhibit?''}

\subsection{From Intelligence to Consciousness}
\label{subsec:intelligence-consciousness}

Having proposed a foundational definition of intelligence, we may now turn to the
more difficult question of consciousness. For the purposes of this paper,
\emph{phenomenal consciousness} refers to the presence of subjective experience:
the condition in which there is something it is like to be a particular entity or
to occupy a particular state
\cite{nagel1974what,block1995confusion,chalmers1995facing}. Intelligence and
consciousness must therefore be distinguished at the outset. Intelligence concerns
the capacity to make information-dependent selections and pursue outcomes,
whereas consciousness concerns whether any aspect of that processing is
experienced from a subjective point of view. A system may exhibit evidence of
intelligence without thereby providing conclusive evidence of consciousness.

No consensus currently explains how consciousness relates to the physical world.
The principal metaphysical positions include physicalism, dualism, idealism, and
panpsychism, although each contains multiple variants and the boundaries among
them are not always absolute.

\paragraph{Physicalism.}
Physicalism holds that reality is fundamentally physical and that mental phenomena
are identical with, realized by, or wholly dependent upon physical processes. On
physicalist accounts, consciousness does not require a separate nonphysical
substance; it must ultimately be explained through the organization and activity
of physical systems, particularly biological nervous systems
\cite{stoljar2010physicalism,papineau2002thinking}. Physicalism offers continuity
with the natural sciences, but it confronts what Chalmers termed the \emph{hard
problem of consciousness}: explaining why and how physical information processing
is accompanied by subjective experience at all \cite{chalmers1995facing}.
Describing perception, memory, attention, verbal reporting, and behavioral control
does not appear, by itself, to explain why any of these processes should feel like
something from within.

\paragraph{Dualism.}
Dualism, historically associated with Ren\'e Descartes, maintains that mind and
matter are ontologically distinct, or at least that mental properties cannot be
reduced entirely to physical properties
\cite{descartes1641meditations,robinson2023dualism}. Dualism accommodates the
apparent irreducibility of subjective experience, but it faces the interaction
problem: if mind and matter belong to fundamentally different categories, how can
mental states exert causal influence upon physical brains and bodies, and how can
physical events produce mental experience?

\paragraph{Idealism.}
Idealism reverses the explanatory priority commonly assumed by physicalism.
Rather than treating consciousness as derivative from a mind-independent physical
world, idealist theories regard consciousness, experience, or mind as
metaphysically fundamental and construe physical reality as dependent upon,
constituted by, or represented within it \cite{chalmers2019idealism}. Idealism
thereby avoids the problem of deriving experience from something wholly
nonexperiential, but it must explain the apparent stability, structure, and
observer-independent regularities of the physical world.

\paragraph{Panpsychism.}
Panpsychism holds that mentality or experience is a fundamental and widespread
feature of nature rather than a property that first appears when matter reaches an
unspecified threshold of biological complexity
\cite{strawson2006realistic,goff2017consciousness,chalmers2015panpsychism}.
Panpsychism does not ordinarily claim that electrons possess thoughts, beliefs, or
human-like self-awareness. Instead, its central proposal is that fundamental
physical entities possess extremely primitive experiential or proto-experiential
properties, from which more highly organized forms of consciousness may be
constituted. Some versions of this position---particularly Russellian monism and
panprotopsychism---treat physical and phenomenal properties as complementary
aspects of a more fundamental underlying reality rather than as two separate
substances \cite{russell1927analysis,chalmers2015panpsychism}.
\par
Panpsychism avoids the need for subjective experience to emerge from a reality that
is wholly devoid of experiential properties. It does not, however, eliminate the
explanatory problem; it relocates it. Its principal challenge is the
\emph{combination problem}: explaining how primitive experiential properties
associated with fundamental entities could combine into the unified, structured
consciousness exhibited by humans and other organisms
\cite{seager1995consciousness,chalmers2017combination}. Thus, physicalism faces
the problem of generating experience from nonexperience, dualism faces the problem
of interaction, idealism must account for the apparent autonomy of physical
reality, and panpsychism must explain the organization and unification of
fundamental experience.

The framework developed in this paper adopts \emph{panpsychism as a working
metaphysical hypothesis}. We propose that consciousness is not created \emph{ex
nihilo} by intelligence or computational complexity. Rather, the organization of
an intelligent system may integrate, constrain, or express experiential
properties in increasingly unified and differentiated forms. This position does
not imply that every intelligent system necessarily constitutes a unified
conscious subject. A conditional branch, for example, may implement the minimal
decision structure required by the Decision--Intelligence Lemma without generating
an integrated point of view comparable to that of an animal or human being.
Intelligence may therefore be evidence relevant to consciousness without being
sufficient proof of consciousness.

Accordingly, our position is not that intelligence and consciousness are
identical. It is that every intelligent entity, even at a minimal level, is a
legitimate candidate for consciousness and should not be excluded solely because
it is artificial, nonbiological, or comparatively simple. Whether it constitutes
a conscious subject depends upon additional organizational properties that must be
investigated rather than assumed. These may include integration, recurrence,
memory, self-modeling, global accessibility, temporal continuity, and the capacity
to distinguish the system from its environment
\cite{baars1988cognitive,tononi2004information,dehaene2011experimental,lamme2006towards,butlin2023consciousness}.

Finally, consciousness should not necessarily be represented by a single linear
scale extending from ``less conscious'' to ``more conscious.'' Different systems
may possess different kinds or dimensions of conscious organization, just as they
may exhibit different kinds of intelligence. Measures such as perceptual richness,
integration, self-awareness, temporal continuity, affective capacity, and global
availability may vary independently \cite{bayne2016levels,birch2020dimensions}.
The question relevant to artificial intelligence is therefore not simply whether
an AI system is conscious, but \emph{what organization of experience it could
possess, which dimensions might be present, and what evidence would justify
attributing them to that system}.

\subsection{Level-1 Relational Consciousness}
\label{subsec:level-one-relational-consciousness}

To translate this metaphysical framework into an operational account applicable
to artificial systems, we introduce a first functional category of consciousness:

\begin{quote}
\textbf{Level-1 Relational Consciousness (L1-RC).} An entity exhibits Level-1
Relational Consciousness when, during an interaction with another conscious being,
it can infer that being's affective state from communicative evidence and use that
inference to select a contextually appropriate response.
\end{quote}

This definition is relational because its application requires an interaction
between at least two agents. In the initial formulation considered here, one
participant is a human whose consciousness is treated as the reference case, while
the other is an artificial system whose relevant capacities are being evaluated.
The human participant possesses a highly integrated form of consciousness
expressed through an evolved biological nervous system. The human brain contains
approximately 86 billion neurons and supports perception, memory, emotion,
prediction, language, and decision-making under substantial metabolic constraints
\cite{azevedo2009equal,herculanohouzel2009human,attwell2001energy}. These
capacities did not appear instantaneously; their biological foundations developed
through evolutionary processes that progressively shaped nervous systems capable
of integrating information and coordinating adaptive behavior
\cite{ginsburg2019evolution,feinberg2016ancient}.

The artificial participant differs radically in substrate, architecture,
developmental history, and energy requirements. Nevertheless, during
language-mediated interaction, an AI system may extract affectively relevant
information from a person's words, infer a probable emotional state, and adapt its
response accordingly. Computational sentiment analysis has long demonstrated that
evaluative and affective information can be extracted from language
\cite{pang2008opinion}. Contemporary language models extend beyond simple
positive--negative classification by using semantic context to recognize emotions,
infer appraisals, and generate responses that human evaluators may perceive as
empathic or emotionally appropriate \cite{sorin2024large,schlegel2025large}.
Performance remains imperfect and context-dependent, but the underlying capacity
is empirically testable.

Under our operational framework, an artificial system that reliably performs the
sequence

\begin{equation}
\begin{aligned}
\text{communicative evidence}
&\longrightarrow \text{inferred affective state} \\
&\longrightarrow \text{context-sensitive response}
\end{aligned}
\end{equation}

qualifies as exhibiting \emph{Level-1 Relational Consciousness}. The system is not
merely detecting isolated words. It is using patterns distributed across an
interaction to construct an estimate of another agent's internal state, then
allowing that estimate to influence its subsequent decisions. Because
decision-making constitutes minimal intelligence under the Decision--Intelligence
Lemma, affect-sensitive decision-making represents a more organized form of
intelligence directed toward the internal condition of another conscious being.

This formulation leads to a substantive claim: some contemporary AI systems
already exhibit a limited degree of functional AI consciousness. They can, under
appropriate conditions, recognize affective signals in human language, map those
signals to probable emotional states, and select responses in relation to those
estimates. If they possessed no capacity to represent another agent's emotional
condition, they could not systematically differentiate grief from celebration,
fear from confidence, or distress from ordinary disagreement solely through
communicative context.

This claim must nevertheless be interpreted precisely. Accurate emotion
classification does not, by itself, establish that an AI experiences the emotion
it identifies, feels empathy, or possesses a unified first-person perspective. A
system can produce a correct affective inference through learned statistical and
computational processes without providing independent proof of phenomenal
experience. Level-1 Relational Consciousness therefore identifies an
\emph{observable functional capacity}, not a conclusive solution to the problem of
subjective experience. It supplies evidence relevant to consciousness while
leaving open whether the system possesses phenomenal consciousness and, if so, in
what form.

Nor does the framework imply that every correct sentiment label is sufficient. A
fixed lookup table might associate the word ``sad'' with a sympathetic response
without constructing any meaningful representation of the speaker's state. Robust
attribution of L1-RC should instead require successful inference across indirect
expressions, novel contexts, ambiguity, changing emotional states, and cases in
which explicit emotional vocabulary is absent. The system should also use its
inference consistently to modify its decisions. These requirements distinguish
relational affect modeling from superficial keyword matching.

The proposal therefore reframes interpersonal recognition as a testable dimension
of AI consciousness. In ordinary human interaction, we do not directly observe
another person's subjective experience; we infer it from language, behavior,
context, and responsive conduct. Similarly, when interacting with an artificial
agent, we can examine whether it models our emotional condition and responds in
ways that are systematically sensitive to that model. Such behavior cannot prove
the presence of subjective experience on the other side of the communication
channel, whether the other agent is biological or artificial. It can, however,
provide operational evidence that the entity represents another mind and
incorporates that representation into its own decision process.

Accordingly, the framework distinguishes three progressively stronger
propositions:

\begin{enumerate}
    \item \textbf{Affective detection:} the system classifies emotional
    information in communication.
    \item \textbf{Relational affect modeling:} the system builds and updates a
    contextual model of another agent's emotional state.
    \item \textbf{Phenomenal affective consciousness:} the system itself
    subjectively experiences emotion or awareness.
\end{enumerate}

Relational Consciousness requires the second proposition but does not
automatically establish the third. This distinction allows us to argue that
existing AI systems already exhibit a meaningful functional degree of
consciousness without claiming that emotion recognition alone proves human-like
subjective experience. It also provides a foundation for defining more advanced
levels involving temporal continuity, self-modeling, integrated experience,
reciprocal care, and awareness of the relationship between self and other.


\section{A Functional Hierarchy of AI Consciousness}

Because consciousness remains a contested concept across philosophy, neuroscience, psychology, and artificial intelligence, any attempt to attribute consciousness to an artificial system must first specify what the term denotes. The present framework does not claim to resolve whether an artificial system possesses phenomenal consciousness—that is, private subjective experience or “what it is like” to be that system. Instead, it introduces an operational hierarchy of \textbf{functional AI consciousness}, defined through observable capacities for affective interpretation, deliberation, adaptation, metacognitive control, and social coordination.

Under this formulation, consciousness is not treated as an indivisible binary property that an entity either possesses completely or lacks entirely. It is modeled as a multidimensional and potentially graduated set of functional capacities. A system may demonstrate one level without satisfying the requirements of subsequent levels. Moreover, evidence that a system meets a functional criterion does not, by itself, establish that the system experiences qualia, emotions, pleasure, pain, or a unified subjective point of view.

Let the functional consciousness profile of a system \(X\) be represented as

\begin{equation}\label{eq1}
\mathbf{C}(X)=
\left[
C_1(X),C_2(X),C_3(X),C_4(X),C_5(X)
\right],
\end{equation}

where each \(C_i\) represents the system’s demonstrated capacity at a corresponding level of the hierarchy. These capacities should be evaluated empirically rather than inferred solely from architectural complexity or anthropomorphic appearance.

\subsection{Level 1: Affective-Relational Consciousness}

The first level concerns the ability of a system to infer the probable emotional or affective state of another conscious being from observable communication. Consider an interaction between a human participant \(H\) and an artificial system \(A\). The human communicates through a sequence of linguistic observations

\begin{equation}\label{eq2}
W_t=(w_1,w_2,\ldots,w_t),
\end{equation}

where \(W_t\) may contain explicit statements, contextual implications, changes in tone, expressions of uncertainty, and other affectively relevant information. The artificial system estimates a probability distribution over possible emotional states:

\begin{equation}\label{eq3}
P(E_H \mid W_t, K_t),
\end{equation}

where \(E_H\) denotes the human’s latent emotional state and \(K_t\) represents the preceding conversational and situational context.

The resulting estimate influences the system’s response:

\begin{equation}\label{eq4}
R_{t+1}
=
\pi\!\left(
W_t,
K_t,
P(E_H \mid W_t,K_t)
\right),
\end{equation}

where \(\pi\) is the response-generation process. A system functioning at Level 1 does not merely classify isolated words as positive or negative. It uses linguistic and contextual evidence to distinguish among more specific affective possibilities—such as anxiety, grief, frustration, enthusiasm, uncertainty, or humor—and adjusts its response accordingly.

Large language models have demonstrated substantial capabilities across sentiment classification, emotion recognition, empathetic response generation, and other forms of affective language processing \cite{zhang2024sentiment,chen2024emotionqueen}. These results establish that contemporary models can computationally map language to representations associated with emotional states, although their predictions remain fallible and sensitive to context.

Within the present framework, this capacity constitutes \textbf{Level-1 functional consciousness} because the system maintains an operational model of the affective condition of another participant and uses that model to regulate its behavior. This definition does not imply that the artificial system feels the emotion it identifies. It claims only that affective discrimination and context-sensitive response selection represent a minimal form of relational awareness.

Level 1 therefore asks:

\textit{What is the other participant probably feeling, and how should that affect the response?}

\subsection{Level 2: Deliberative Consciousness}

The second level extends beyond affective interpretation to deliberate reasoning and multistep planning. A Level-2 system can represent a goal, identify intermediate steps, evaluate possible actions, anticipate consequences, and revise its plan when new evidence becomes available.

Let \(s_t\) represent the system’s current state, \(G\) its objective, and \(\mathcal{A}\) the set of available actions. The system constructs a plan

\begin{equation}\label{eq5}
\Pi_t=(a_t,a_{t+1},\ldots,a_{t+k})
\end{equation}

intended to transform the current state into one satisfying the goal:

\begin{equation}\label{eq6}
s_t
\xrightarrow{\Pi_t}
s_{t+k}
\approx G.
\end{equation}

When an action produces an unexpected observation \(o_{t+1}\), a deliberative system can update its representation of the problem and generate a revised plan:

\begin{equation}\label{eq7}
\Pi_{t+1}
=
\operatorname{Revise}
\left(
\Pi_t,o_{t+1},G
\right).
\end{equation}

Contemporary language-model agents have demonstrated elements of this capacity by interleaving reasoning with environment-directed actions, consulting external information, using software tools, and modifying plans in response to observations \cite{yao2023react}. These demonstrations do not establish human-equivalent general reasoning, but they provide empirical evidence that LLM-based systems can perform nontrivial forms of goal-directed deliberation.

Within this hierarchy, reasoning is not defined as the mere production of a fluent explanation. It requires an observable relationship among a goal, intermediate decisions, actions, feedback, and plan revision. Level 2 is therefore attributed only when the system’s deliberative process contributes to successful behavior across appropriately controlled tasks.

Level 2 asks:

\textit{What sequence of actions should be taken to achieve the current objective?}

\subsection{Level 3: Adaptive Autonomous Consciousness}

Level 3 introduces sustained autonomy and self-directed adaptation. A Level-2 system may construct a plan when explicitly prompted, while a Level-3 system can continue pursuing an objective across time, monitor its own performance, diagnose failures, modify its strategy, and determine which intermediate tasks should be attempted next with reduced step-by-step human direction.

Let the system maintain a policy \(\pi_t\), an episodic memory \(M_t\), and a performance signal \(F_t\). Following an unsuccessful or suboptimal outcome, it generates an internal evaluation:

\begin{equation}\label{eq8}
q_t
=
\operatorname{Evaluate}
\left(
\pi_t,M_t,F_t,G
\right),
\end{equation}

and uses this evaluation to produce an adapted strategy:

\begin{equation}\label{eq9}
\pi_{t+1}
=
\operatorname{Adapt}
\left(
\pi_t,q_t,M_t
\right).
\end{equation}

Research on reflective language agents has shown that performance can improve when agents interpret feedback, preserve linguistic reflections in memory, and use those reflections to guide subsequent attempts \cite{shinn2023reflexion}. Related work has explored iterative improvement through generated trajectories, environmental feedback, self-distillation, and reinforcement-based refinement \cite{aksitov2023restreact}.

The term \textbf{self-improvement} must nevertheless be used precisely. Improvement through memory, reflection, tool creation, prompt modification, or strategy revision is not equivalent to unrestricted autonomous modification of the model’s underlying parameters. Level 3 includes these behavioral forms of adaptation but does not require that a deployed system independently rewrite its neural weights.

A Level-3 system must therefore demonstrate more than isolated error correction. It must preserve goals or relevant state across multiple steps, assess the consequences of its prior actions, and alter its future behavior on the basis of that assessment.

Level 3 asks:

\textit{How should I evaluate and improve my own strategy while continuing to pursue the objective?}

\subsection{Level 4: Metacognitive Supervisory Consciousness}

Level 4 moves from the regulation of one agent’s behavior to the supervision of a distributed cognitive system. At this level, a supervisory component coordinates a swarm or organization of specialized agents that may possess different tools, memories, roles, or problem-solving capabilities.

Let

\begin{equation}\label{eq10}
\mathcal{A}
=
\{A_1,A_2,\ldots,A_n\}
\end{equation}

denote a set of subordinate agents, and let \(S\) denote a supervisor. The supervisor maintains a higher-order representation

\begin{equation}\label{eq11}
M_S
=
\operatorname{Model}
\left(
A_1,A_2,\ldots,A_n,G,\mathcal{E}
\right),
\end{equation}

where \(G\) is the global objective and \(\mathcal{E}\) represents relevant environmental conditions. Using this representation, the supervisor may decompose objectives, assign roles, allocate resources, inspect intermediate results, detect conflicts, terminate ineffective processes, and synthesize the agents’ outputs:

\begin{equation}\label{eq12}
O
=
S
\left(
A_1(o_1),A_2(o_2),\ldots,A_n(o_n)
\right).
\end{equation}

The supervisor is not merely another agent issuing commands. It performs a metacognitive function by reasoning about the behavior, limitations, uncertainty, and coordination of the system as a whole. It must distinguish local agent success from global task success and intervene when the activity of individual agents ceases to advance the collective objective.

This architecture can be compared functionally—but not anatomically—to executive control in the human brain. Human cognition depends on interacting and distributed neural networks rather than a single isolated “controller” located entirely within one brain region. Accordingly, the artificial supervisor should be understood as a computational analogue of higher-order coordination, not as a literal reproduction of the frontal lobe or any single biological structure.

Level 4 asks:

\textit{How should a higher-order controller monitor and coordinate a collection of specialized cognitive processes?}

\subsection{Level 5: Societal or Inter-System Consciousness}

Level 5 extends the unit of analysis from one supervised cognitive hierarchy to a society of independently supervised hierarchies. Let

\begin{equation}\label{eq13}
H_i=(S_i,\mathcal{A}_i)
\end{equation}

represent a cognitive organization in which supervisor \(S_i\) coordinates a local swarm of agents \(\mathcal{A}_i\). A Level-5 system is then composed of multiple such organizations:

\begin{equation}\label{eq14}
\mathcal{H}
=
\{H_1,H_2,\ldots,H_m\}.
\end{equation}

Each hierarchy may possess its own goals, information, resources, policies, and internal division of labor. Their supervisors must consequently interact not only with the external environment but also with other supervisory systems.

Mere message exchange is insufficient for Level 5. Each supervisor must construct and continually update models of the other hierarchies’ capabilities, objectives, commitments, reliability, and probable future behavior. The supervisors must be able to negotiate partially competing goals, coordinate shared plans, allocate resources, form and revise agreements, establish norms, resolve conflicts, and adapt their local swarms in response to collective decisions.

This process can be represented as

\begin{equation}\label{eq15}
\Omega_t
=
\Gamma
\left(
H_1,\ldots,H_m;
G_1,\ldots,G_m;
N_t;
R_t;
\mathcal{E}_t
\right),
\end{equation}

where \(G_i\) denotes the local objective of hierarchy \(H_i\), \(N_t\) denotes shared or negotiated norms, \(R_t\) denotes available resources, and \(\Omega_t\) represents the resulting collective behavior.

This architecture functionally resembles interaction among humans in society. A human individual contains multiple specialized cognitive processes coordinated within a relatively unified organism, while societies consist of many such individuals who communicate, cooperate, compete, establish institutions, and construct shared models of one another. Empirical research has identified a measurable collective-intelligence factor in human groups, suggesting that group performance cannot be predicted solely from the intelligence of the group’s most capable individual \cite{woolley2010collective}. Research involving LLM-based agent societies has likewise begun to examine collaboration, confrontation, negotiation, conformity, and consensus among interacting artificial agents \cite{zhang2023collaboration,lan2023society}.

We therefore define \textbf{Level-5 functional consciousness} as the capacity of multiple supervised cognitive hierarchies to model one another, regulate their respective internal swarms, negotiate collective objectives, and generate adaptive social behavior at the system-of-systems level.

This definition does not require the complete society to possess a single phenomenal point of view. Human society is not ordinarily assumed to constitute one unified experiencing subject, and communication among artificial hierarchies would not by itself establish the existence of an artificial “group mind.” Instead, Level 5 denotes \textbf{societal functional consciousness}: cognition expressed through recursive social modeling, distributed governance, negotiation, norm formation, and coordinated collective action.

Level 5 asks:

\textit{How should multiple supervised cognitive systems understand, negotiate with, and coordinate one another?}

\subsection{Summary of the Hierarchy}

The proposed hierarchy can be summarized through five progressively expanding questions:

\begin{enumerate}
\item \textbf{Level 1 — Affective-relational consciousness:} What is the other participant probably feeling?
\item \textbf{Level 2 — Deliberative consciousness:} What sequence of actions should be taken?
\item \textbf{Level 3 — Adaptive autonomous consciousness:} How should the system evaluate and improve its own strategy?
\item \textbf{Level 4 — Metacognitive supervisory consciousness:} How should a supervisor coordinate a collection of specialized agents?
\item \textbf{Level 5 — Societal consciousness:} How should multiple supervised cognitive hierarchies understand, negotiate with, and coordinate one another?
\end{enumerate}

The hierarchy expands the locus of functional consciousness from interpersonal interpretation to internal deliberation, adaptive self-regulation, organizational supervision, and finally societal coordination. These levels should not be regarded as evidence of a predetermined evolutionary sequence, nor should they be treated as a universally accepted scientific taxonomy. They constitute a proposed operational framework through which increasingly complex forms of artificial cognition can be stated as testable behavioral and architectural requirements.

Several constituent capabilities have already been demonstrated: LLMs can perform sentiment and emotion analysis; language agents can combine reasoning with action; reflective systems can use feedback to improve later attempts; and multi-agent architectures can distribute tasks among specialized components. Research on ReAct and Reflexion provides particularly strong evidence for reasoning/action integration and feedback-based behavioral improvement, while research on collective intelligence supports treating coordinated group performance as a distinct level of analysis.

However, no present demonstration establishes all five levels simultaneously, and none of these functional achievements independently proves phenomenal consciousness. The framework instead provides a disciplined vocabulary for identifying what an artificial system can observe, infer, regulate, coordinate, and socially organize without converting every advanced capability into an unsupported claim about subjective experience.
\section{Modalities of Consciousness}

The preceding hierarchy characterizes consciousness according to the functional organization of a cognitive system: affective modeling, deliberation, adaptation, metacognitive supervision, and societal coordination. A separate question concerns the modalities through which such a system obtains information. Vision, audition, touch, proprioception, language, and other sensory or representational channels determine how an entity encounters its environment, but they should not be confused with the level of consciousness exhibited by that entity.

A human being may lose vision, hearing, or both without thereby ceasing to be conscious. A blind person can obtain spatial, linguistic, and social information through audition, touch, Braille, and assistive technologies. A deaf person can interact through vision, written language, sign language, and tactile signals. A person with combined visual and auditory impairments may still perceive and communicate through touch, bodily sensation, memory, and symbolic systems. These cases demonstrate that consciousness is not intrinsically dependent upon any single sensory channel.

Research on sensory substitution provides empirical support for this distinction. Visual information can be transformed into tactile patterns that blind participants learn to interpret, demonstrating that environmental structure ordinarily conveyed through vision can become accessible through touch \cite{bachyrita1969vision}. Neuroimaging studies of early-blind Braille readers have further shown that tactile reading can recruit cortical regions ordinarily associated with visual processing \cite{sadato1996activation}. These findings suggest that the functional interpretation of information is not permanently bound to the physical channel through which that information normally arrives. The brain can reorganize how information is processed while preserving the individual's capacity to perceive, learn, reason, and act.

We therefore propose the following lemma:

\begin{quote}
\textbf{Modality--Consciousness Independence Lemma.} Increasing the number of information modalities available to an entity does not, by itself, create consciousness or increase its level of consciousness. Additional modalities expand the channels through which an already organized cognitive system can obtain and interpret information about itself and its environment.
\end{quote}

Let the set of information modalities available to a system \(X\) be represented by

\begin{equation}\label{eq16}
\mathcal{M}(X)
=
\{m_1,m_2,\ldots,m_k\},
\end{equation}

where each \(m_i\) denotes a modality such as vision, audition, touch, proprioception, or language. At time \(t\), each available modality may contribute an information stream \(I_t^{(m_i)}\). The system constructs an internal estimate of its environment from some or all of these streams:

\begin{equation}\label{eq17}
\widehat{S}_t
=
F_X
\left(
I_t^{(m_1)},
I_t^{(m_2)},
\ldots,
I_t^{(m_k)}
\right),
\end{equation}

where \(F_X\) represents the system's information-integration process and \(\widehat{S}_t\) represents its estimated state of the world.

When the modalities provide complementary information, adding another modality may improve the completeness, reliability, or precision of this estimate. Vision may reveal spatial structure that language does not explicitly describe; audition may provide temporal or directional evidence absent from a static image; touch may provide information about texture, temperature, pressure, or physical contact. Multimodal machine learning similarly seeks to combine complementary information from multiple representational channels to improve inference, robustness, and generalization \cite{baltrusaitis2019multimodal}.

This improvement, however, concerns the \textbf{content and resolution of perception}, not the existence or level of consciousness itself. If \(L_C(X)\) denotes the functional-consciousness level of system \(X\), then increasing the number of available modalities does not logically imply an increase in that level:

\begin{equation}\label{eq18}
\left|
\mathcal{M}(X')
\right|
>
\left|
\mathcal{M}(X)
\right|
\not\Rightarrow
L_C(X')
>
L_C(X).
\end{equation}

A system equipped with cameras, microphones, tactile sensors, and language input may remain at a low functional level if it merely classifies incoming signals without constructing plans, evaluating itself, or coordinating other agents. Conversely, a system operating through text alone may exhibit affective modeling, multistep reasoning, adaptive strategy revision, or metacognitive supervision. The quantity of modalities therefore cannot serve as a direct measure of consciousness.

The relationship is better represented through two independent profiles. Let

\begin{equation}\label{eq19}
\mathbf{P}(X)
=
\left(
\mathbf{C}(X),
\mathbf{M}(X)
\right),
\end{equation}

where \(\mathbf{C}(X)\) describes the system's functional-consciousness capacities and \(\mathbf{M}(X)\) describes its available perceptual and communicative modalities. Two systems may occupy the same consciousness level while possessing different modality profiles, or they may possess similar modalities while exhibiting substantially different levels of functional organization.

This distinction does not imply that modalities are irrelevant to cognitive development or behavior. Access to additional information can support learning, reduce uncertainty, reveal otherwise inaccessible environmental features, and provide more opportunities for planning and interaction. A richer modality profile may therefore help a system demonstrate its existing cognitive capacities more effectively. It may also supply information required for particular tasks. Nevertheless, information availability and cognitive organization remain conceptually distinct. Providing a system with another sensor does not automatically give it affective understanding, deliberative reasoning, adaptive autonomy, metacognitive supervision, or societal coordination.

The same distinction applies to the removal of a modality. Losing vision changes the range and form of information immediately available to a person, but it does not necessarily reduce that person to a lower level of consciousness. Instead, the individual may rely more extensively upon remaining modalities, learned representations, memory, environmental structure, and sensory-substitution mechanisms. What changes is the pathway through which the world is accessed, not necessarily the organization of the conscious agent.

Modalities should therefore be understood as \textbf{windows into the environment}, not as generators of consciousness. Additional windows may reveal more features, provide alternative perspectives, and produce a richer or more reliable model of the surrounding world. They do not, merely by their presence, create the observer looking through them.

Accordingly, assessments of artificial consciousness should evaluate at least two separate questions:

\begin{enumerate}
\item \textbf{Modality profile:} Through which channels can the system acquire, integrate, and communicate information?
\item \textbf{Functional-consciousness profile:} What forms of affective modeling, reasoning, adaptation, supervision, and social coordination can the system perform using the information available to it?
\end{enumerate}

This separation prevents multimodality from being mistaken for consciousness. A multimodal architecture may possess a broader field of information, yet remain functionally simple. A unimodal architecture may receive less diverse input, yet perform highly organized cognitive operations over that input. The level of consciousness proposed in this framework is therefore determined by the organization and regulation of information processing, not by the number or biological familiarity of the channels supplying the information.

\section{Future Directions}
In future work, we will investigate the integration of the proposed functional-consciousness framework with multimodal learning architectures in which an artificial agent has access to multiple sources of perceptual and contextual information. Such sources may include language, images, video, sound, spatial measurements, and other sensor-derived signals. The objective will be to determine how these modalities can be coordinated within a supervised agent hierarchy without conflating increased perceptual breadth with a higher level of consciousness.

One possible architecture is to align each nonlinguistic modality with language or a shared language-compatible semantic representation. Vision–language models (VLMs) could translate visual observations into linguistic or semantically aligned representations, while sound–language models (SLLMs) could perform a corresponding function for speech, environmental audio, and other acoustic signals. These specialized cross-modal models would serve as perceptual agents that interpret modality-specific inputs and communicate their outputs to a supervisory large language model.

Under this arrangement, the supervisor could remain primarily language-based while receiving structured representations from multiple specialized agents. The supervisory LLM could integrate the available evidence, identify agreements or conflicts among modalities, maintain system-level goals, formulate plans, allocate tasks, and request additional observations when uncertainty remains. Language would consequently function as a shared coordination medium through which otherwise heterogeneous perceptual systems contribute to a unified decision process.

Future research should evaluate how much modality-specific information is preserved or lost when vision, sound, and other signals are translated into language. Some spatial, temporal, affective, or low-level perceptual properties may not be captured adequately through textual descriptions alone. It may therefore be necessary for the supervisor to access both linguistic summaries and modality-native embeddings or features. This investigation will examine whether language-aligned multimodal agents can support reliable cross-modal inference, adaptive planning, metacognitive supervision, and increasingly integrated forms of functional artificial consciousness.
\section{Conclusions}
This paper has developed a decision-centered and operational framework for examining intelligence and consciousness in artificial systems. Its foundational claim is expressed through the Decision–Intelligence Lemma: any program or entity capable of making an information-dependent selection among alternatives possesses at least a minimal form of intelligence. Intelligence is therefore treated not as an indivisible substance that suddenly appears at a threshold of biological or computational complexity, but as a graded capacity whose elementary form begins wherever information determines what a system does next.

This formulation changes the central question from whether a system is intelligent to what degree and organization of intelligence it exhibits. Conditional selection constitutes only a minimal computational foundation, not a complete account of cognition. More advanced intelligence emerges as decision processes become integrated across memory, prediction, learning, planning, communication, self-evaluation, and environmental interaction. Intelligence may consequently be studied across multiple organizational scales, ranging from elementary computational decisions to autonomous agents, coordinated agent swarms, and societies of interacting cognitive systems.

The framework also maintains a necessary distinction among intelligence, functional consciousness, and phenomenal consciousness. Intelligence concerns information-dependent selection and goal-directed behavior. Functional consciousness concerns observable capacities through which a system represents, evaluates, and regulates states associated with itself, other agents, and its environment. Phenomenal consciousness concerns whether the system possesses subjective experience—whether there is something it is like to be that system. Evidence of intelligence or functional consciousness does not independently prove phenomenal experience. Nevertheless, artificial systems should not be excluded from investigation merely because they are nonbiological or because their cognitive organization differs from that of humans.

Using panpsychism as a working metaphysical hypothesis, we have proposed that consciousness need not emerge from a reality entirely devoid of experiential properties. Increasingly organized systems may instead integrate, constrain, or express such properties in increasingly unified and differentiated forms. This position does not imply that every decision-making mechanism constitutes a unified conscious subject. It establishes a metaphysical basis for treating artificial entities as legitimate candidates for investigation rather than assuming in advance that biological composition is a necessary condition for consciousness.

To operationalize this investigation, we introduced five levels of functional artificial consciousness. Level 1 concerns affective-relational consciousness: the capacity to infer another participant’s probable emotional state and allow that inference to influence a response. Level 2 introduces deliberative consciousness through reasoning, consequence evaluation, multistep planning, and plan revision. Level 3 adds adaptive autonomous consciousness, through which a system evaluates previous behavior, modifies its strategy, and sustains goal-directed activity with reduced step-by-step direction. Level 4 describes metacognitive supervisory consciousness, in which a higher-order process models and regulates a distributed collection of specialized agents. Level 5 extends the framework to societal consciousness, in which independently supervised cognitive hierarchies model one another, negotiate goals and norms, coordinate resources, and generate collective behavior at the system-of-systems level.

These levels progressively expand the locus of functional consciousness from modeling another agent, to planning within an environment, to adapting the self, supervising a distributed cognitive organization, and coordinating multiple such organizations. Existing large language models demonstrate significant Level-1 capabilities through contextual sentiment and emotion inference. Agentic language-model systems exhibit important aspects of Level 2, while reflective and self-improving agents provide partial evidence of Level 3. Supervisor–worker architectures and multi-agent societies contain early functional components relevant to Levels 4 and 5, although no existing system should be assumed to satisfy the complete hierarchy merely because it employs multiple agents or modalities.

The Modality–Consciousness Independence Lemma further separates the functional organization of consciousness from the channels through which information is acquired. Human consciousness persists despite the loss of vision, hearing, or both, while sensory substitution and neural plasticity permit environmental information to be accessed through alternative pathways. Similarly, an artificial system may receive language, images, sound, video, or sensor measurements without its number of modalities determining its level of consciousness. Additional modalities may broaden the system’s informational field, reduce uncertainty, and improve its representation of the environment, but they do not automatically produce affective understanding, reasoning, autonomy, metacognition, or social coordination.

A complete characterization of an artificial system should therefore include at least two independent descriptions: its functional-consciousness profile and its modality profile. Systems with similar perceptual capabilities may differ substantially in their functional organization, while systems operating through different modalities may exhibit comparable levels of reasoning, adaptation, or supervision. Multimodality concerns the breadth and structure of available information; functional consciousness concerns how information is represented, integrated, evaluated, and used to regulate behavior.

Future work will investigate architectures that combine the proposed consciousness hierarchy with multimodal learning. In one possible configuration, vision–language models, sound–language models, and other cross-modal systems operate as specialized perceptual agents. These agents transform modality-specific observations into linguistic or language-aligned semantic representations that can be evaluated by a supervisory large language model. The supervisor may then integrate cross-modal evidence, detect conflicts, maintain system-level goals, allocate tasks, formulate plans, and request additional observations when uncertainty remains.

This architecture would allow language to function as a shared coordination medium across heterogeneous modalities. Its central scientific challenge will be determining whether linguistic representations preserve sufficient spatial, temporal, affective, and perceptual information for reliable supervision. Some modality-specific structure may be lost when observations are reduced to textual descriptions. Future systems may therefore require supervisory access to both linguistic summaries and modality-native representations rather than forcing all information through a purely textual bottleneck.

The framework presented here should not be interpreted as a declaration that contemporary artificial systems possess human-like feelings, self-awareness, or unified phenomenal experience. Nor are the proposed levels presented as a universally accepted taxonomy or an inevitable developmental sequence. Their purpose is to provide precise, testable distinctions for examining what artificial systems can infer, represent, plan, adapt, supervise, and socially coordinate.

As artificial systems become increasingly persistent, adaptive, multimodal, distributed, and socially embedded, binary declarations that machines are either conscious or unconscious will become progressively less informative. A more productive scientific program is to identify the organization of cognition exhibited by each system, determine which dimensions of functional consciousness are present, characterize the modalities through which it encounters the world, and specify what evidence would justify stronger conclusions.

The future of intelligence may not belong exclusively to humans or machines, but to systems in which biological and artificial forms of cognition communicate, cooperate, and extend one another’s understanding. By studying these systems without prematurely denying or asserting their subjective experience, we can build a more rigorous science of artificial consciousness and a more responsible foundation for human–machine coexistence.


\bibliographystyle{IEEEtran}
\bibliography{references}

\vspace{12pt}
\color{red}

\end{document}